\documentclass[11pt]{article}
\usepackage[margin=1in]{geometry}
\usepackage{amsmath,amssymb,amsthm}
\usepackage{xurl}
\usepackage{hyperref}
\sloppy
\let\oldus\_ \renewcommand{\_}{\oldus\allowbreak}
\newcommand{\Icc}{\operatorname{Icc}}
\title{Conjecture 00000008538: the radical and quotient of an incidence algebra \\ (false without finiteness; true for finite lattices)}
\date{}
\begin{document}
\maketitle

\section*{Statement and verdict}
\textbf{Conjecture.} ``A lattice algebra is the incidence algebra of a lattice. The Jacobson radical of an
incidence algebra is its strictly lower triangular part; and its quotient is the Mobius algebra, of dimension
equal to the number of lattice elements.''

\textbf{Verdict: DISPROVED as written.} The statement has no finiteness hypothesis. Take $K=\mathbb Q$ and the infinite,
locally finite lattice $L=\mathbb N^{op}$ (the naturals with the reversed order, which has a top element). Then
the conjunction fails: if the radical were the strictly triangular part $S$, the quotient $I/J$ would be $K$-linearly $K^L$, of
dimension $|K^{\mathbb N}|=2^{\aleph_0}$, not $|L|=\aleph_0$. The statement is TRUE for finite lattices (second theorem below).
So the honest boundary is: it holds for finite lattices and fails for infinite locally finite lattices.

\section*{Reading (all choices explicit)}
\begin{itemize}
\item $K$ is a field, $L$ a lattice with a top element (so the Mobius algebra is unital) such that $L$ is locally finite,
which Mathlib needs to make $I(L,K)$ a ring. The incidence algebra is Mathlib's \texttt{IncidenceAlgebra K L}: functions $f(x,y)$ vanishing unless
$x\le y$, with $(fg)(x,y)=\sum_{x\le z\le y}f(x,z)g(z,y)$. This is the upper-triangular convention; ``lower'' is the transposed
convention (passing to $L^{op}$). The counterexample does not depend on this, because the strictly triangular part is defined by $f(x,x)=0$.
\item The Jacobson radical is \texttt{Ring.jacobson} of the ring $I(L,K)$, the intersection of the maximal left ideals. The strictly triangular part is
\texttt{StrictlyTriangular K L} $=\{f: f(x,y)\neq 0\Rightarrow x<y\}$, compared with the radical as sets.
\item The Mobius algebra is the monoid algebra $K[(L,\wedge,\top)]$ (\texttt{MobiusAlgebra K L}); ``its quotient is'' is a $K$-algebra isomorphism of $I/J$ with it.
``Dimension equal to the number of lattice elements'' is \texttt{Module.rank K (I/J) = Cardinal.mk L}, which is the right reading for infinite $L$.
\end{itemize}
\textbf{Not refuted / contested.} A reader may take ``the number of lattice elements'' to presuppose that $L$ is finite. Under that reading the
statement is the finite theorem, which we prove (it holds). We refute the statement as written, with no finiteness assumption. We do not decide
whether the isomorphism clause alone fails for $\mathbb N^{op}$ (only the radical-plus-rank conjunction is refuted).

\section*{Disproof}
Let $I=I(L,K)$, $J=\operatorname{Rad}(I)$, $S=\{f: f(x,x)=0\ \forall x\}$ (equal to the strictly triangular set: $f(x,y)\ne0$ forces $x\le y$, and
$f(x,x)=0$ excludes $x=y$; Lean \texttt{strict\_iff}). Suppose $J=S$ as sets.
The diagonal map $\delta:I\to K^L$, $\delta(f)(a)=f(a,a)$, is a $K$-algebra homomorphism (\texttt{diag}). It is surjective: for $g\in K^L$ the
function $(a,b)\mapsto[a=b]g(a)$ is an incidence function (\texttt{diag\_surj}). Its kernel is $S=J$. Hence $I/J\cong K^L$ as $K$-modules
(\texttt{quotLinEquiv}: restriction of scalars of the quotient, \texttt{quotEquivOfEq}, first isomorphism theorem), so
$\operatorname{rank}_K(I/J)=\operatorname{rank}_K(K^L)$. For infinite $L$, the Erdos--Kaplansky theorem in Mathlib (\texttt{rank\_fun\_infinite}) gives $\operatorname{rank}_K K^L=|K^L|$
(\texttt{rank\_ne}). For $K=\mathbb Q$, $L=\mathbb N^{op}$: $|\mathbb Q^{\mathbb N}|=\aleph_0^{\aleph_0}=2^{\aleph_0}>\aleph_0=|L|$. So the third clause
fails, contradicting the conjecture (\texttt{main}). Thus no field and infinite lattice pair with $J=S$ can satisfy the rank clause, and
$\mathbb Q$, $\mathbb N^{op}$ (locally finite, with top, instances from Mathlib) is one.

\section*{Boundary: the finite theorem}
\texttt{finite\_main : FiniteConjecture}: for every field $K$ and finite lattice $L$, (1) $J=S$, (2) $I/J\cong K[(L,\wedge)]$ as $K$-algebras, (3) $\dim_K I/J=|L|$.
Proof sketch (as in the earlier package). $J\subseteq S$: $f\mapsto f(a,a)$ is a surjective ring map onto the field $K$, whose radical is $0$.
$S\subseteq J$: elements of $S$ are nilpotent (induction on $|\Icc(a,b)|$), so $1-rf$ is a unit for every $r$, hence $f$ lies in every maximal left ideal.
The diagonal map gives $I/J\cong K^L$, and $K^L\cong K[L,\wedge]$ by Solomon's idempotents $e_a=\sum_{y\le a}\mu(y,a)y$ using Mathlib's
\texttt{mu} (\texttt{mobPhi\_surj}) plus equality of dimensions.

\section*{Lean}
File \texttt{lean/Conjecture8538/Basic.lean} (identical to \texttt{Tlmc/C8538.lean}), imports only Mathlib. The statement block defines \texttt{MeetMonoid},
\texttt{MobiusAlgebra}, \texttt{StrictlyTriangular}, \texttt{Conjecture} (all fields, all locally finite lattices with a top), \texttt{FiniteConjecture}
and \texttt{Claim := $\neg$ Conjecture}. Main theorem \texttt{C8538.main : C8538.Claim}; boundary theorem \texttt{C8538.finite\_main : C8538.FiniteConjecture}.
\texttt{\#print axioms} for both reports only \texttt{propext}, \texttt{Classical.choice}, \texttt{Quot.sound}; no \texttt{sorry}, no \texttt{native\_decide}.
Sources: Mathlib (\texttt{IncidenceAlgebra}, \texttt{rank\_fun\_infinite}, \texttt{Cardinal.aleph0\_power\_aleph0}); no external literature is cited.
\end{document}
